\section{Quarkus - Grundlagen und Einsatz in der Anwendung}
\setauthor{MT}

Quarkus ist ein modernes, quelloffenes Java-Framework, das speziell für den Einsatz 
in containerisierten und Cloud-nativen Umgebungen entwickelt wurde. Ziel des Frameworks 
ist es, den Prozess der Bereitstellung von Java-Applikationen in Cloud-Umgebungen zu 
vereinfachen und zu optimieren. Dabei unterstützt Quarkus Entwicklerinnen und Entwickler 
dabei, Applikationen auch ohne tiefgehende Kenntnisse der zugrundeliegenden Infrastruktur 
bereitzustellen.

Ein wesentliches Merkmal von Quarkus ist seine Kompatibilität mit einer Vielzahl 
etablierter Frameworks und Standards. Dazu zählen unter anderem Eclipse MicroProfile, 
RESTEasy (JAX-RS), Hibernate ORM (JPA), Apache Kafka sowie Spring. Diese breite 
Unterstützung ermöglicht es, bestehende Kenntnisse und Bibliotheken direkt einzusetzen, 
ohne grundlegende Umstrukturierungen vornehmen zu müssen.

Quarkus hebt sich durch seine Leichtgewichtigkeit von anderen Java-Frameworks ab. 
Konkret äußert sich dies in einer geringen Startzeit, einem niedrigen Speicherverbrauch 
sowie der einfachen Möglichkeit, Projekte in Cloud-Umgebungen bereitzustellen.


\section{Grundprinzip von Quarkus}
\setauthor{MT}

Ein zentrales Merkmal von Quarkus liegt in der Art, wie Java-Anwendungen kompiliert 
und ausgeführt werden. Um dies zu verstehen, ist es notwendig, die zwei grundlegenden 
Kompilierungsansätze zu unterscheiden: die Just-In-Time-Kompilierung (JIT) und die 
Ahead-of-Time-Kompilierung (AOT).

\subsection{Just-In-Time-Kompilierung (JIT)}

Im klassischen Java-Ausführungsmodell wird Quellcode zunächst vom Java-Compiler in 
plattformunabhängigen Bytecode übersetzt. Dieser Bytecode wird anschließend von der 
Java Virtual Machine (JVM) interpretiert. Die JIT-Kompilierung ist dabei ein 
Optimierungsmechanismus der JVM: Häufig ausgeführte Codeabschnitte werden zur Laufzeit 
in nativen Maschinencode übersetzt, was die Ausführungsgeschwindigkeit erhöht.

Der Nachteil dieses Ansatzes liegt darin, dass die JVM beim Start der Applikation 
zunächst hochgefahren werden muss und die JIT-Optimierungen erst schrittweise während 
der Laufzeit greifen. Dies führt zu längeren Startzeiten und einem erhöhten 
Speicherverbrauch – Faktoren, die in containerisierten Umgebungen mit häufigen 
Neustarts besonders nachteilig sind.

\subsection{Ahead-of-Time-Kompilierung (AOT)}

Die Ahead-of-Time-Kompilierung verfolgt einen grundlegend anderen Ansatz: Der gesamte 
Anwendungscode wird bereits vor der Ausführung, also zum Zeitpunkt des Build-Prozesses, 
in nativen Maschinencode übersetzt. Quarkus nutzt hierfür GraalVM Native Image, ein 
Werkzeug, das Java-Anwendungen in plattformspezifische, ausführbare Binärdateien 
kompiliert.

Das Ergebnis ist eine native ausführbare Datei, die ohne eine laufende JVM auskommt. 
Dies hat folgende Vorteile:

\begin{itemize}
    \item \textbf{Schnelle Startzeit:} Da keine JVM initialisiert werden muss und alle 
    Optimierungen bereits zur Build-Zeit vorgenommen wurden, startet die Anwendung 
    in Millisekunden.
    \item \textbf{Geringer Speicherverbrauch:} Der Wegfall der JVM-Laufzeitumgebung 
    reduziert den RAM-Bedarf erheblich.
    \item \textbf{Eignung für containerisierte Umgebungen:} Kurze Startzeiten und 
    geringer Ressourcenverbrauch machen AOT-kompilierte Anwendungen besonders 
    geeignet für Kubernetes und ähnliche Plattformen.
\end{itemize}

Als Nachteil ist die deutlich längere Build-Zeit im Vergleich zur JIT-Kompilierung 
zu nennen, da die gesamte Kompilierung und Optimierung vorab stattfindet.

\subsection{Hot Reload}

Für den Entwicklungsprozess bietet Quarkus einen sogenannten Hot-Reload-Mechanismus, 
bekannt als \textit{Quarkus Dev Mode}. Dabei werden Änderungen am Quellcode automatisch 
erkannt und ohne Neustart der Anwendung übernommen. Dies beschleunigt den 
Entwicklungszyklus erheblich, da zeitaufwändige Neustarts entfallen. Der Dev Mode 
verwendet intern JIT-Kompilierung, um diese schnelle Rückkopplungsschleife zu 
ermöglichen – die AOT-Kompilierung wird lediglich für den finalen Build eingesetzt.


\section{Quarkus in der vorliegenden Applikation}
\setauthor{MT}

Quarkus wird in der vorliegenden Anwendung als Hauptframework für das Backend verwendet. 
Für die Datenbankanbindung kommt Hibernate ORM mit Panache zum Einsatz, das eine 
vereinfachte Abstraktion des JPA-Standards bietet und den Datenbankzugriff mit weniger 
Boilerplate-Code ermöglicht.

Die Anwendung wird mittels AOT-Kompilierung als native ausführbare Datei gebaut. 
Obwohl die Build-Zeit dadurch im Vergleich zur klassischen JIT-Kompilierung länger 
ausfällt, überwiegen die Vorteile im Betrieb: Die Applikation weist eine deutlich 
verkürzte Startzeit auf, benötigt weniger Arbeitsspeicher und ist optimal für den 
Einsatz in einer containerisierten Umgebung geeignet. Diese Eigenschaften sind 
insbesondere in Cloud-Deployments von Vorteil, wo Ressourceneffizienz und 
Skalierbarkeit zentrale Anforderungen darstellen.
